% Options for packages loaded elsewhere
\PassOptionsToPackage{unicode}{hyperref}
\PassOptionsToPackage{hyphens}{url}
\documentclass[
  ignorenonframetext,
]{beamer}
\newif\ifbibliography
\usepackage{pgfpages}
\setbeamertemplate{caption}[numbered]
\setbeamertemplate{caption label separator}{: }
\setbeamercolor{caption name}{fg=normal text.fg}
\beamertemplatenavigationsymbolsempty
% remove section numbering
\setbeamertemplate{part page}{
  \centering
  \begin{beamercolorbox}[sep=16pt,center]{part title}
    \usebeamerfont{part title}\insertpart\par
  \end{beamercolorbox}
}
\setbeamertemplate{section page}{
  \centering
  \begin{beamercolorbox}[sep=12pt,center]{section title}
    \usebeamerfont{section title}\insertsection\par
  \end{beamercolorbox}
}
\setbeamertemplate{subsection page}{
  \centering
  \begin{beamercolorbox}[sep=8pt,center]{subsection title}
    \usebeamerfont{subsection title}\insertsubsection\par
  \end{beamercolorbox}
}
% Prevent slide breaks in the middle of a paragraph
\widowpenalties 1 10000
\raggedbottom
\AtBeginPart{
  \frame{\partpage}
}
\AtBeginSection{
  \ifbibliography
  \else
    \frame{\sectionpage}
  \fi
}
\AtBeginSubsection{
  \frame{\subsectionpage}
}
\usepackage{iftex}
\ifPDFTeX
  \usepackage[T1]{fontenc}
  \usepackage[utf8]{inputenc}
  \usepackage{textcomp} % provide euro and other symbols
\else % if luatex or xetex
  \usepackage{unicode-math} % this also loads fontspec
  \defaultfontfeatures{Scale=MatchLowercase}
  \defaultfontfeatures[\rmfamily]{Ligatures=TeX,Scale=1}
\fi
\usepackage{lmodern}
\ifPDFTeX\else
  % xetex/luatex font selection
\fi
% Use upquote if available, for straight quotes in verbatim environments
\IfFileExists{upquote.sty}{\usepackage{upquote}}{}
\IfFileExists{microtype.sty}{% use microtype if available
  \usepackage[]{microtype}
  \UseMicrotypeSet[protrusion]{basicmath} % disable protrusion for tt fonts
}{}
\makeatletter
\@ifundefined{KOMAClassName}{% if non-KOMA class
  \IfFileExists{parskip.sty}{%
    \usepackage{parskip}
  }{% else
    \setlength{\parindent}{0pt}
    \setlength{\parskip}{6pt plus 2pt minus 1pt}}
}{% if KOMA class
  \KOMAoptions{parskip=half}}
\makeatother
\setlength{\emergencystretch}{3em} % prevent overfull lines
\providecommand{\tightlist}{%
  \setlength{\itemsep}{0pt}\setlength{\parskip}{0pt}}
% Auto-generated by infrastructure/rendering/_pdf_latex_helpers.py::extract_math_font_preamble.
% Minimal header passed to Pandoc via -H so Beamer slides pick up the same math font
% as the combined-PDF path. Adjust the manuscript's preamble.md to change the font globally.
\usepackage{unicode-math}
\setmathfont{latinmodern-math.otf}

% Slide layout defaults for warning-clean scientific decks.
\usepackage{etoolbox}
\IfFileExists{xurl.sty}{\usepackage{xurl}}{}
\IfFileExists{seqsplit.sty}{\usepackage{seqsplit}}{\newcommand{\seqsplit}[1]{#1}}
\protected\def\breakseq#1{\seqsplit{#1}}
\protected\def\breaktt#1{\begingroup\ttfamily\seqsplit{#1}\endgroup}
\setlength{\emergencystretch}{6em}
\tolerance=5000
\hbadness=10000
\hfuzz=1pt
\setlength{\tabcolsep}{2pt}
\AtBeginEnvironment{longtable}{\tiny\renewcommand{\arraystretch}{0.86}\setlength{\tabcolsep}{1pt}}
\AtBeginEnvironment{tabular}{\tiny\renewcommand{\arraystretch}{0.86}\setlength{\tabcolsep}{1pt}}

% Natbib and cross-reference fallbacks — slides don't load natbib
% or cleveref, but combined-PDF manuscript prose may emit these
% commands. The fallback renders citations as a bracketed key list
% and cross-references as detokenized labels so slides don't fail on
% undefined control sequences or raw underscores. \providecommand is
% a no-op if packages load later.
\providecommand{\citep}[1]{[#1]}
\providecommand{\citet}[1]{#1}
\providecommand{\citealp}[1]{#1}
\providecommand{\citeauthor}[1]{#1}
\providecommand{\citeyear}[1]{#1}
\providecommand{\cref}[1]{\texttt{\detokenize{#1}}}
\providecommand{\Cref}[1]{\texttt{\detokenize{#1}}}

\newtheorem{proposition}{Proposition}
\newtheorem{hypothesis}{Hypothesis}
\usepackage{bookmark}
\IfFileExists{xurl.sty}{\usepackage{xurl}}{} % add URL line breaks if available
\urlstyle{same}
% fallback for those not using the hyperref driver hyperxmp:
\makeatletter
\@ifundefined{xmpquote}{\newcommand{\xmpquote}[1]{#1}}{}
\makeatother
\hypersetup{
  hidelinks,
  pdfcreator={LaTeX via pandoc}}

\author{\texorpdfstring{}{}}
\date{}

\begin{document}

\section{Scope, Related Work, and Positioning}\label{sec:scope}

\begin{frame}[allowframebreaks]{Scope, Related Work, and Positioning}
This section situates the exemplar scientifically and states explicit
boundaries. The goal is not to compete with monographs on nonlinear
programming \breakseq{[@nocedal2006numerical; @bertsekas1999nonlinear]}, but to
show how a minimal, test-backed optimization story fits the template's
reproducibility and rendering stack \breakseq{[@peng2011reproducible]}.
\end{frame}

\begin{frame}[allowframebreaks]{Classical gradient methods}
\protect\phantomsection\label{classical-gradient-methods}
Smooth unconstrained minimization via first-order updates has a long
lineage, from Cauchy's early descent perspective
\breakseq{[@cauchy1847methode]} to modern treatments of convex problems
\breakseq{[@boyd2004convex]} and accelerated gradient schemes
\breakseq{[@nesterov2013gradient]}. Polyak's classical discussion of gradient
convergence factors remains relevant when interpreting empirical
iteration counts \breakseq{[@polyak1964some]}. The present manuscript restricts
attention to \textbf{fixed-step} gradient descent on a \textbf{convex
quadratic}, where rates reduce to scalar linear recurrences in the error
(\breakseq{[@sec:results]}, \breakseq{[@eq:scalar\_linear\_update]}).
\end{frame}

\begin{frame}[fragile,allowframebreaks]{Adaptive and stochastic
extensions}
\protect\phantomsection\label{adaptive-and-stochastic-extensions}
Practical machine-learning optimizers (e.g., Adam \breakseq{[@kingma2014adam]})
introduce momentum, adaptive preconditioning, or noise from
minibatching. Those methods are \textbf{out of scope} for
\breaktt{template\_code\_project}: the exemplar deliberately keeps the
algorithm minimal so that failures (divergent \(\alpha\), iteration
caps) are interpretable without confounding from stochastic sampling or
line-search logic.
\end{frame}

\begin{frame}[fragile,allowframebreaks]{What this project proves about
the template}
\protect\phantomsection\label{what-this-project-proves-about-the-template}
The scientific claims through \breakseq{[@sec:introduction]},
\breakseq{[@sec:methodology]}, and \breakseq{[@sec:results]} are standard textbook
material. The \textbf{non-standard} contribution is procedural:
configuration in \breaktt{manuscript/config.yaml} drives
\breaktt{run\_convergence\_experiment()}, figures, CSV exports, and
\texttt{\{\{RESULT\_*\}\}} substitution
(\breaktt{scripts/z\_generate\_manuscript\_variables.py}) so that PDF,
HTML, and validation logs refer to the same numbers. That pattern is
what downstream projects should copy---whether the domain is
optimization, differential equations, or Bayesian workflows.
\end{frame}

\begin{frame}[fragile,allowframebreaks]{Explicit limitations}
\protect\phantomsection\label{explicit-limitations}
\begin{enumerate}
\tightlist
\item
  \textbf{Dimensionality}: Default experiments emphasize \(d = 1\) with
  \(A = I\) for transparent plotting; the benchmark figure explores
  \(d > 1\) only with identity Hessians, so no ill-conditioning effects
  appear.
\item
  \textbf{Step-size policy}: Only constant \(\alpha\) is implemented in
  \breaktt{src/optimizer.py}; there is no Wolfe or Armijo backtracking.
\item
  \textbf{Global optimization}: Convexity is assumed; no basin-hopping
  or restarts are studied.
\item
  \textbf{Numerical model}: Double-precision floating point only; no
  interval or arbitrary-precision analysis.
\end{enumerate}

These limitations are intentional: they narrow the failure surface so
that infrastructure concerns---tests, logging, figure registration, and
PDF cross-references---remain visible rather than buried under
algorithmic complexity.
\end{frame}

\end{document}
